% Синтетический кейс: нумерация как в LaTeX (общие счетчики, \numberwithin, \tag,
% subequations, безномерные окружения). Реальных статей в корпусе нет.
\documentclass[12pt]{article}
\usepackage[T2A]{fontenc}
\usepackage[utf8]{inputenc}
\usepackage[russian]{babel}
\usepackage{amsmath,amssymb,amsthm}

\theoremstyle{plain}
\newtheorem{theorem}{Теорема}[section]
\newtheorem{proposition}[theorem]{Предложение}
\theoremstyle{definition}
\newtheorem{definition}{Определение}
\newtheorem*{notation}{Обозначение}

\numberwithin{equation}{section}

\begin{document}

\section{Первый раздел}

\begin{definition}
Первое определение со своим счетчиком.
\end{definition}

\begin{theorem}
Утверждение без метки.
\end{theorem}

\begin{equation}
  x = y
\end{equation}

\section{Второй раздел}

\begin{definition}\label{def:two}
Второе определение.
\end{definition}

\begin{theorem}\label{thm:two}
Теорема с меткой.
\end{theorem}

\begin{proposition}
Предложение делит счетчик с теоремой.
\end{proposition}

\begin{notation}
Обозначение без номера.
\end{notation}

\begin{align}
  a &= b, \\
  c &= d.
\end{align}

\begin{equation}\label{eq:tagged}
  e = f \tag{A}
\end{equation}

\begin{equation}
  g = h \tag*{B}
\end{equation}

\begin{equation*}
  i = j
\end{equation*}

\begin{subequations}\label{eq:sub}
\begin{align}
  k &= l, \label{eq:sub-a}\\
  m &= n. \label{eq:sub-b}
\end{align}
\end{subequations}

\begin{multline}\label{eq:multi}
  o + p + q \\
  + r + s = t.
\end{multline}

Ссылки: определение~\ref{def:two}, теорема~\ref{thm:two}, формулы~\eqref{eq:tagged},
\eqref{eq:sub}, \eqref{eq:sub-a}, \eqref{eq:sub-b} и~\eqref{eq:multi}.

\end{document}
